% ===========================================================================
% Relatório do projeto — Arquitetura e Desenho de Software (ADS)
% Sistema Monitor: requisitos de qualidade, decisões arquiteturais e avaliação
%
% Compilar:   make            (pdflatex + bibtex + pdflatex x2)
% Dados:      make dados      (copia as tabelas geradas pelas medições; ver README.md)
%
% Antes de entregar: preencher os campos a vermelho (\porescrever) e os autores.
% ===========================================================================
\documentclass[11pt, a4paper]{article}

% ----------------------------------------------------- dados da capa -------
\newcommand{\Mestrado}{Engenharia Informática}
\newcommand{\UnidadeCurricular}{Arquitetura e Desenho de Software}
\newcommand{\Autores}{%
  João Afonso, N.º 97990\\
  Rodrigo Gorjão, N.º 82972}
\newcommand{\Docente}{Vitor Basto-Fernandes}
\newcommand{\Titulo}{Monitor: Arquitetura Orientada a Atributos de Qualidade}
\newcommand{\Subtitulo}{Métricas de Qualidade, Escalabilidade e Manutenção de um Sistema de Monitorização Distribuído}
\newcommand{\Entrega}{Projeto --- entrega final}
\newcommand{\Data}{9 de dezembro de 2026}
\newcommand{\TituloCurto}{Monitor: arquitetura orientada a QA}
\newcommand{\Organizacao}{https://github.com/ads-iscte-rg-ja}
\newcommand{\Repositorio}{\Organizacao/monitor-platform}
% Nome de um repositório da organização, com ligação.
\newcommand{\repo}[1]{\href{\Organizacao/#1}{\texttt{#1}}}

% ---------------------------------------------------------------------------
% Preâmbulo: pacotes, aspeto e macros do relatório.
% ---------------------------------------------------------------------------

\usepackage[T1]{fontenc}
\usepackage[utf8]{inputenc}
\usepackage[portuguese]{babel}
\usepackage{lmodern}
\usepackage{microtype}
\usepackage[a4paper, margin=2.5cm, headheight=14pt]{geometry}
\usepackage{setspace}
\onehalfspacing

% Matemática
\usepackage{amsmath, amssymb, mathtools, bm}

% Tabelas
\usepackage{array, booktabs, tabularx, longtable, multirow, makecell}
\usepackage{colortbl}
\usepackage{threeparttable}
\usepackage{pdflscape}

% Figuras
\usepackage{graphicx}
\usepackage{float}
\usepackage[section]{placeins}  % figuras e tabelas não passam para a secção seguinte
\usepackage[font=small, labelfont=bf, labelsep=period]{caption}
\usepackage{subcaption}
\usepackage{tikz}
\usetikzlibrary{arrows.meta, positioning, shapes.geometric, shapes.multipart,
                fit, calc, backgrounds, decorations.pathreplacing}
\usepackage{pgfplots}
\pgfplotsset{compat=1.16}

% Diversos
\usepackage{xcolor}
\usepackage[inline]{enumitem}
\usepackage{pifont}
\usepackage{csquotes}
\usepackage{url}
\usepackage{fancyhdr}
\usepackage{titlesec}
\usepackage[hidelinks, unicode, pdfusetitle]{hyperref}

% ---------------------------------------------------------------------------
% Cores (inspiradas na identidade do Iscte; ajustáveis)
% ---------------------------------------------------------------------------
\definecolor{iscte}{RGB}{0, 51, 102}
\definecolor{acento}{RGB}{0, 112, 192}
\definecolor{cinza}{RGB}{110, 110, 110}
\definecolor{fundotab}{RGB}{235, 241, 248}
\definecolor{bom}{RGB}{0, 128, 0}
\definecolor{mau}{RGB}{190, 30, 45}
\definecolor{parcial}{RGB}{200, 130, 0}

% ---------------------------------------------------------------------------
% Notação UML e BPMN (estilos TikZ partilhados pelas figuras)
% ---------------------------------------------------------------------------
\tikzset{
  % UML: pacote (pasta com separador); o nome vai no corpo
  umlpkg/.style={draw=iscte, fill=fundotab, minimum height=7mm, minimum width=22mm, font=\small\ttfamily,
    align=center, append after command={\pgfextra{\draw[iscte, fill=fundotab]
      ([yshift=-0.1pt]\tikzlastnode.north west) rectangle ++(6mm,1.6mm);}}},
  % UML: pacote que contém outros (nome no separador)
  umlpkgbox/.style={draw=iscte, rounded corners=0pt, inner sep=7pt, append after command={\pgfextra{
      \node[draw=iscte, fill=white, anchor=south west, font=\small\bfseries\ttfamily, inner sep=2pt,
        minimum height=4mm] at ([yshift=-0.1pt]\tikzlastnode.north west) {#1};}}},
  % UML: componente (ícone no canto superior direito)
  umlcomp/.style={draw=iscte, fill=fundotab, minimum height=9mm, minimum width=22mm, font=\small, align=center,
    inner sep=3pt, append after command={\pgfextra{
      \draw[iscte, fill=white] ([xshift=-4.2mm, yshift=-1.2mm]\tikzlastnode.north east) rectangle ++(3mm,-2.6mm);
      \draw[iscte, fill=white] ([xshift=-5mm, yshift=-1.7mm]\tikzlastnode.north east) rectangle ++(1.6mm,-0.6mm);
      \draw[iscte, fill=white] ([xshift=-5mm, yshift=-2.9mm]\tikzlastnode.north east) rectangle ++(1.6mm,-0.6mm);}}},
  % UML: artefacto/base de dados (estereótipo no texto)
  umlart/.style={draw=iscte, fill=white, minimum height=8mm, minimum width=20mm, font=\small, align=center},
  % UML: nó de implantação (caixa 3D); o nome é dado por label
  umlnode/.style={draw=cinza, inner sep=7pt, append after command={\pgfextra{
      \draw[cinza] (\tikzlastnode.north west) -- ++(2mm,2mm) -| ([xshift=2mm, yshift=2mm]\tikzlastnode.south east)
        -- (\tikzlastnode.south east);
      \draw[cinza] (\tikzlastnode.north east) -- ++(2mm,2mm);}}},
  % UML: classe ou interface com compartimentos
  umlclass/.style={draw=iscte, fill=fundotab, rectangle split, rectangle split parts=2, font=\scriptsize\ttfamily,
    align=left, inner sep=3pt},
  % UML: setas
  umldep/.style={-{Straight Barb[length=2mm]}, dashed, thick, iscte},           % dependência («use»)
  umlreal/.style={-{Triangle[open, length=2.5mm, width=2.5mm]}, dashed, thick, iscte}, % realização
  umlassoc/.style={-{Straight Barb[length=2mm]}, thick, iscte},                % associação navegável / mensagem
  umlmsg/.style={-{Stealth[length=2mm]}, thick, iscte},                        % mensagem síncrona
  umlret/.style={-{Straight Barb[length=2mm]}, dashed, iscte},                  % resposta
  umlconn/.style={thick, iscte},                                               % conector/caminho de comunicação
  % UML (atividade) e BPMN
  bpmntask/.style={draw=iscte, fill=fundotab, rounded corners=3pt, minimum height=9mm, minimum width=19mm,
    text width=18mm, font=\scriptsize, align=center, inner sep=2pt},
  bpmnstart/.style={circle, draw=iscte, thick, fill=white, minimum size=5mm, inner sep=0pt},
  bpmnend/.style={circle, draw=iscte, line width=1.8pt, fill=white, minimum size=5mm, inner sep=0pt},
  bpmngw/.style={diamond, draw=iscte, thick, fill=white, minimum size=7mm, inner sep=0pt, font=\small\bfseries},
  bpmnflow/.style={-{Stealth[length=2mm]}, thick, iscte},
  actstart/.style={circle, fill=iscte, minimum size=3.5mm, inner sep=0pt},
  actend/.style={circle, draw=iscte, thick, fill=white, minimum size=4.5mm, inner sep=0pt,
    path picture={\fill[iscte] (path picture bounding box.center) circle (1.2mm);}},
}

% ---------------------------------------------------------------------------
% Títulos de secção e cabeçalhos
% ---------------------------------------------------------------------------
\titleformat{\section}{\Large\bfseries\color{iscte}}{\thesection}{0.8em}{}
\titleformat{\subsection}{\large\bfseries\color{iscte}}{\thesubsection}{0.8em}{}
\titleformat{\subsubsection}{\normalsize\bfseries\color{iscte}}{\thesubsubsection}{0.8em}{}

\pagestyle{fancy}
\fancyhf{}
\fancyhead[L]{\small\color{cinza}\TituloCurto}
\fancyhead[R]{\small\color{cinza}ADS --- Projeto}
\fancyfoot[C]{\small\thepage}
\renewcommand{\headrulewidth}{0.4pt}

% Nomes em português de Portugal
\addto\captionsportuguese{%
  \renewcommand{\contentsname}{Índice}%
  \renewcommand{\listfigurename}{Lista de Figuras}%
  \renewcommand{\listtablename}{Lista de Tabelas}%
  \renewcommand{\refname}{Referências}%
  \renewcommand{\abstractname}{Resumo}%
  \renewcommand{\tablename}{Tabela}%
  \renewcommand{\figurename}{Figura}%
  \renewcommand{\appendixname}{Anexo}%
}

% ---------------------------------------------------------------------------
% Macros
% ---------------------------------------------------------------------------
% Marcas usadas nas tabelas comparativas
\newcommand{\Sim}{\textcolor{bom}{\ding{51}}}          % abordado
\newcommand{\Parc}{\textcolor{parcial}{\ding{108}}}    % abordado parcialmente
\newcommand{\Nr}{\textcolor{cinza}{--}}                % não reportado
\newcommand{\Nao}{\textcolor{mau}{\ding{55}}}          % ausente por desenho


% Colunas de tabela
\newcolumntype{L}[1]{>{\raggedright\arraybackslash}p{#1}}
\newcolumntype{C}[1]{>{\centering\arraybackslash}p{#1}}
\newcolumntype{Y}{>{\raggedright\arraybackslash}X}

% Caixas de definição/observação
\newcommand{\caixa}[2]{%
  \begin{center}\begin{tikzpicture}
    \node[draw=acento, fill=fundotab, rounded corners=2pt, inner sep=8pt,
          text width=0.9\linewidth, align=justify] {\textbf{#1.}\ #2};
  \end{tikzpicture}\end{center}}

% ---------------------------------------------------------------------------
% Macros específicas do relatório de ADS
% ---------------------------------------------------------------------------
% Texto ainda por escrever: aparece a vermelho para não passar despercebido.
\newcommand{\porescrever}[1]{\textcolor{mau}{\itshape [Por escrever: #1]}}

% Inclui uma tabela/figura gerada a partir das medições (make dados).
% Se o ficheiro ainda não existir, mostra uma caixa com o nome esperado.
\newcommand{\gerado}[1]{%
  \IfFileExists{gerado/#1}{\input{gerado/#1}}%
    {\caixa{Dados em falta}{\texttt{gerado/\detokenize{#1}} --- correr \texttt{make dados}.}}}

% Identificadores de cenários, decisões e lacunas (ligam ao plano e aos ADRs).
\newcommand{\cen}[1]{\textsf{#1}}
\newcommand{\dec}[1]{\textsf{D#1}}


\hypersetup{pdftitle={\Titulo: \Subtitulo}, pdfsubject={ADS -- Projeto}}

% Folga extra para evitar linhas a sair da margem (termos em \texttt não hifenizam)
\setlength{\emergencystretch}{2em}

% Páginas só com tabelas/figuras começam no topo
\makeatletter
\setlength{\@fptop}{0pt}
\setlength{\@fpbot}{0pt plus 1fil}
\makeatother

\begin{document}

% ------------------------------------------------------------- capa -------
\pagenumbering{Alph}
\begin{titlepage}
  \thispagestyle{empty}
  \centering
  \IfFileExists{figuras/iscte-logo.png}%
    {\includegraphics[width=0.62\textwidth]{figuras/iscte-logo.png}}%
    {\begin{tikzpicture}
       \node[draw=iscte, line width=1pt, minimum width=6cm, minimum height=2.2cm,
             text=iscte, font=\Large\bfseries] {iscte};
     \end{tikzpicture}}
  \par\vspace{2.2cm}

  {\Large Mestrado em \Mestrado\par}
  \vspace{0.6cm}
  {\LARGE\bfseries\color{iscte} \UnidadeCurricular\par}
  \vspace{1.6cm}

  {\Large\bfseries \Titulo\par}
  \vspace{0.4cm}
  {\large \Subtitulo\par}
  \vspace{0.6cm}
  {\normalsize\itshape \Entrega\par}
  \vspace{1.6cm}

  \begin{flushleft}
    \large
    \textbf{Autores:}\\[0.2cm]
    \Autores\\[0.8cm]
    \textbf{Docente:} \Docente\\[0.8cm]
    \textbf{Repositórios:} \url{\Organizacao}\\[0.1cm]
    {\normalsize \repo{monitor-central}, \repo{monitor-remote}, \repo{monitor-go-sdk},
    \repo{monitor-hello-world}, \repo{monitor-api}, \repo{monitor-platform}}\\[0.4cm]
    \textbf{Data:} \Data
  \end{flushleft}
  \vfill
\end{titlepage}

% ---------------------------------------------------- pré-textuais --------
\pagenumbering{roman}
\begin{abstract}
O Monitor é um sistema de monitorização de máquinas e aplicações: um agente em cada máquina recolhe métricas
e erros, guarda-os localmente e entrega-os a um servidor central, que avalia alertas e os mostra num
\emph{dashboard}. O sistema funcionava, mas a arquitetura
estava implícita no código e não se sabia quanto aguentava. Este trabalho torna-a explícita e melhora-a de
forma orientada a atributos de qualidade: cenários com medidas verificáveis (escalabilidade para 1\,000
\emph{hosts}, disponibilidade, manutenção, segurança), decisões tomadas com o ADD, fundamentadas em táticas e
registadas em ADR, e cada iteração medida antes e depois, com métricas de código e testes de carga. As medições mudaram o desenho. A \emph{baseline} mostrou que, depois de uma hora de indisponibilidade
do central, a recuperação deixava o \emph{dashboard} inutilizável (p95 de pelo menos 10~s); com escritas em lote
e controlo de admissão, o p95 desceu para 262~ms e o \emph{backlog} passou a ser absorvido duas vezes mais
depressa, sem perdas nem duplicados. A medição refutou também o risco de vizinhos ruidosos, cuja proteção não foi
implementada, e mostrou que, com várias instâncias do central, o estrangulamento passou para a base de dados: o
relatório regista o resultado e o gatilho para a escalar, em vez de ajustar o alvo.

\porescrever{rever o resumo no fim, com os resultados de \dec{7}, \dec{13} e \dec{19}.}

\medskip
\noindent\textbf{Palavras-chave:} arquitetura de software, atributos de qualidade, escalabilidade,
métricas de qualidade, táticas arquiteturais, ADR.
\end{abstract}


\clearpage
\tableofcontents
\clearpage

% --------------------------------------------------------- texto ----------
\pagenumbering{arabic}
\section{Introdução}
\label{sec:introducao}

\subsection{Contexto e motivação}
\porescrever{o sistema Monitor (central, remote, SDK); projeto proposto pela equipa; relação com EUVI (UI/UX).}

\subsection{Objetivos}
\porescrever{objetivos de negócio G1--G6 (plano §1) e as três prioridades de ADS: métricas de qualidade, escalabilidade e manutenção.}

\subsection{Contribuições e estrutura do relatório}
\porescrever{lista de contribuições e guia das secções.}

\section{Enquadramento teórico}
\label{sec:enquadramento}

Este capítulo resume apenas os conceitos que o relatório usa, e como os usa. A referência principal são
os diapositivos da unidade curricular~\cite{bastofernandes2026}, baseados em Bass et
al.~\cite{bass2012}.

\subsection{Arquitetura de software, estruturas e vistas}

A arquitetura de software de um sistema é o conjunto de estruturas necessárias para raciocinar sobre
ele: elementos de software, as relações entre eles e as propriedades de ambos~\cite{bass2012}. Há três
categorias de estruturas~\cite{bastofernandes2026}:
\begin{itemize}
  \item \textbf{de módulos}: unidades de implementação e as relações de decomposição, uso e camadas.
        Determinam a modificabilidade, porque definem até onde uma alteração se propaga;
  \item \textbf{de componentes e conectores} (C\&C): elementos em execução e a forma como interagem.
        Permitem raciocinar sobre desempenho, disponibilidade e concorrência;
  \item \textbf{de alocação}: a correspondência entre o software e o ambiente, isto é, máquinas,
        ficheiros e equipas. A estrutura de \emph{deployment} serve para raciocinar sobre desempenho,
        integridade dos dados, segurança e disponibilidade.
\end{itemize}
Uma vista é a representação documentada de uma estrutura. O capítulo~\ref{sec:sistema} documenta uma
vista de cada categoria. Segue-se o princípio de só documentar o que traz retorno: as partes que um
recém-chegado precisa de aprender, as que têm risco e as que mudam com frequência (YAGNI).

\subsection{Atributos de qualidade e cenários}

Os sistemas são muitas vezes redesenhados não por lhes faltar funcionalidade, mas por serem difíceis de
manter, escalar ou proteger~\cite{bastofernandes2026}. Os atributos de qualidade (QA) dividem-se em atributos
\textbf{de execução} (disponibilidade, desempenho, usabilidade) e \textbf{de desenvolvimento}
(modificabilidade, testabilidade). Os nomes seguem os diapositivos e o modelo de qualidade da ISO/IEC
25010~\cite{iso25010}. Os QA
não se atingem isoladamente: melhorar um afeta outros, e desenhar é, em parte, escolher
\emph{trade-offs}. Por isso o capítulo~\ref{sec:requisitos} os torna explícitos.

Um requisito de qualidade tem de ser inequívoco e testável. Especifica-se como um cenário de seis partes:
\textbf{fonte} do estímulo, \textbf{estímulo}, \textbf{artefacto} afetado, \textbf{ambiente},
\textbf{resposta} e \textbf{medida da resposta}. A medida é o que torna o cenário verificável, e é ela que
os testes de carga e as métricas de código deste trabalho medem.

Três definições dos diapositivos 03 são usadas na avaliação. Quando as responsabilidades de dois módulos
se sobrepõem, uma única alteração pode afetar ambos; o \textbf{acoplamento} é a probabilidade de isso
acontecer, e é mau para a modificabilidade. A \textbf{coesão} é a
probabilidade de uma alteração que afeta uma responsabilidade de um módulo afetar também as outras, e é
boa quando é alta. A \textbf{testabilidade} é a facilidade com que os testes revelam as faltas que
existem.

\subsection{Táticas e padrões}

Uma tática é uma decisão de desenho que influencia a resposta de um único atributo de qualidade, sem ter
em conta os outros. Os \emph{trade-offs} entre táticas têm de ser controlados explicitamente pelo
arquiteto~\cite{bastofernandes2026}. Os padrões arquiteturais (Layered, Client-Server, Shared-Data,
Broker, MVC, Publish-Subscribe) já incluem os seus \emph{trade-offs}. O Broker, por exemplo, traz
indireção e latência, e torna-se um ponto de ataque, de falha e de estrangulamento. Ao nível do desenho
detalhado usam-se os padrões GoF~\cite{gamma1994}. Neste trabalho os catálogos de táticas por QA
(disponibilidade, modificabilidade, desempenho, segurança, testabilidade) servem duas vezes: para
identificar o que o sistema já faz (secção~\ref{sec:tacticas-existentes}) e para justificar cada decisão
nova.

A análise de um QA pode fazer-se por estimativas \emph{back-of-the-envelope}, por experiências mentais
ou por \emph{checklists}, como na segurança. A análise feita mais tarde no ciclo de vida, sobre o sistema
real, dá mais confiança, mas custa mais~\cite{bastofernandes2026}. É essa a relação entre as estimativas
do capítulo~\ref{sec:requisitos}, baratas e feitas cedo, que apontam os prováveis estrangulamentos, e os
testes de carga, que os confirmam ou desmentem.

\subsection{Desenho orientado a atributos (ADD) e registo de decisões}

Os requisitos arquiteturalmente significativos (ASR) são os que alteram profundamente a arquitetura; são
muitas vezes requisitos de qualidade e derivam frequentemente dos objetivos da própria organização. Uma
\textbf{utility tree} organiza-os: a raiz é a ``utilidade'' do sistema, o primeiro nível são os QA, o
segundo são refinamentos, e as folhas são cenários. Cada folha é classificada pelo
\textbf{valor de negócio} e pelo \textbf{impacto arquitetural}, numa escala de H, M e L. Os cenários (H,H)
são os primeiros a tratar no desenho~\cite{bastofernandes2026}.

O \emph{Attribute-Driven Design} (ADD) é iterativo. Em cada iteração escolhe-se um elemento, reúnem-se
os ASR que lhe dizem respeito, gera-se e testa-se uma solução, e escolhe-se o que tratar na iteração
seguinte. O desenho avança por \emph{generate and test}: cada hipótese é testada e a seguinte parte do
resultado. Neste trabalho o ``teste'' de cada iteração é uma medição: a mesma medição de baseline,
repetida depois de cada alteração.

Cada decisão é registada num \emph{Architecture Decision Record} (ADR)~\cite{nygard2011}, com o contexto,
a decisão e as consequências, a que se acrescentam a teoria que a fundamenta, as alternativas rejeitadas e
a forma de a verificar. Os ADR estão no anexo~\ref{anx:adrs}. Não é uma prática só académica: o
\emph{Technology Radar} da Thoughtworks recomenda-os para adoção desde 2017~\cite{thoughtworks2017adr}, a AWS
publica um guia sobre o processo~\cite{kunce2022adr} e o \emph{Azure Well-Architected Framework} da Microsoft
trata-os como um dos principais entregáveis do arquiteto~\cite{microsoftadr}. O motivo é preservar o porquê de cada
decisão, que o código não mostra, e evitar voltar a discutir o que já foi decidido.

\subsection{Métricas de qualidade de software}

As métricas são escolhidas pelo método \emph{Goal--Question--Metric}~\cite{basili1994}: um objetivo gera
perguntas, e só se mede o que responde a uma delas. As métricas de código usadas são:
\begin{itemize}
  \item \textbf{complexidade ciclomática}~\cite{mccabe1976}: o número de caminhos linearmente
        independentes de uma função, que é também o número de testes necessários para percorrer uma
        base desses caminhos. Mede o esforço de teste;
  \item \textbf{complexidade cognitiva}~\cite{campbell2018}: penaliza o aninhamento e as quebras de
        fluxo, e não a mera sequência de decisões. Mede o esforço de leitura;
  \item \textbf{métricas de pacote de Martin}~\cite{martin2002}: acoplamento aferente (Ca, quem depende do
        pacote) e eferente (Ce, de quem o pacote depende), instabilidade $I = Ce/(Ca+Ce)$, abstração $A$
        (fração de tipos abstratos) e distância à sequência principal $D = |A + I - 1|$. Um pacote estável
        deve ser abstrato, e um instável deve ser concreto; $D$ próximo de 0 indica esse equilíbrio.
\end{itemize}
A aplicação destas métricas ao Monitor, com os limites escolhidos, está na
secção~\ref{sec:gqm}.

\section{O sistema Monitor: arquitetura atual}
\label{sec:sistema}

\subsection{Vista de módulos}
\porescrever{diagrama TikZ da decomposição (plano §5.1) e regras de camadas verificadas por \texttt{depguard}.}

\subsection{Vista de componentes e conectores}
\porescrever{componentes de execução, conectores e fluxo de telemetria.}

\subsection{Vista de alocação}
\porescrever{deployment: host monitorizado, servidor central, serviços externos.}

\subsection{Padrões e tácticas já presentes}
\porescrever{tabela de padrões e tácticas existentes (plano §5.2--5.3).}

\section{Requisitos arquiteturais}
\label{sec:requisitos}

\subsection{Priorização dos atributos de qualidade}
\porescrever{ranking e justificação (plano §2.1).}

\subsection{Trade-offs explícitos}
\porescrever{tabela de conflitos e tácticas de mitigação (plano §2.2).}

\subsection{Utility tree}
\porescrever{árvore com classificação (valor, impacto) (plano §3.1).}

\subsection{Cenários de qualidade}
\porescrever{cenários detalhados SC1--SC5, MO1, AV1, AV4 (plano §3.2).}

\section{Métricas de qualidade e ferramentas}
\label{sec:metricas}

\subsection{Seleção das métricas (GQM)}
\porescrever{tabela objetivo--pergunta--métrica--limite (plano §4.1) e separação entre métricas principais e gates.}

\subsection{Ferramentas e alternativas}
\porescrever{golangci-lint, SonarCloud, govulncheck, k6, Prometheus; alternativas rejeitadas (plano §4.2).}

\subsection{Pipeline de integração contínua}
\porescrever{monitor-platform, workflow reutilizável, gates no código novo (ADR-0002).}

\subsection{Baseline}
\gerado{baseline-lint.tex}
\gerado{baseline-cobertura.tex}

\section{Decisões arquiteturais}
\label{sec:decisoes}

\subsection{Escalabilidade}
\porescrever{\dec{4}--\dec{10}, \dec{19}: problema, tática, alternativas, trade-offs (resumo dos ADRs).}

\subsection{Manutenção e testabilidade}

\subsubsection{Contrato de ingestão explícito em gRPC (\dec{11}, ADR-0003)}
\label{sec:d11}

\paragraph{Problema.}
O contrato entre o remote e o central era implícito (lacuna M2). O remote decidia se punha um registo em
quarentena comparando o \textbf{texto em inglês} da mensagem de erro do central com uma lista fixa de
frases: mudar uma mensagem no central mudava o comportamento do remote sem que nenhum teste falhasse.
As estruturas JSON estavam escritas à mão nos dois repositórios, e a classificação das respostas
(quarentena, dividir o lote, tentar mais tarde, pausar) dependia de pares código HTTP e campo
\texttt{error} documentados apenas no código. É o acoplamento semântico que o cenário \cen{MO2} proíbe.

\paragraph{Decisão.}
O contrato passou a ser um serviço gRPC definido em Protocol Buffers, num repositório próprio,
\texttt{monitor-api}, versionado por \emph{tags} (\texttt{v1.0.0}). O central (servidor) e o remote
(cliente) dependem dele pelo \texttt{go.mod}. Os erros levam um código gRPC e, quando um único registo é
inválido, um detalhe com o identificador do registo e um motivo estável (\texttt{RejectionReason}); o
texto da mensagem deixou de fazer parte do contrato. A Tabela~\ref{tab:erros-grpc} mostra a
correspondência, que preserva os quatro comportamentos do \emph{store-and-forward}. Os endpoints HTTP de
ingestão foram removidos; o \emph{dashboard} continua em HTTP.

\begin{table}[htbp]
  \centering\small
  \caption{Erros do contrato de ingestão e reação do remote}
  \label{tab:erros-grpc}
  \begin{tabularx}{\linewidth}{L{0.3\linewidth}L{0.3\linewidth}Y}
    \toprule
    Situação & Código e detalhe & O remote \\
    \midrule
    Um registo inválido & \texttt{INVALID\_ARGUMENT} + \texttt{RecordRejection} & põe esse registo em quarentena e reenvia os restantes \\
    Lote mal formado & \texttt{INVALID\_ARGUMENT}, sem detalhe & pausa (violação de contrato) \\
    Mensagem acima de 1~MiB & \texttt{RESOURCE\_EXHAUSTED}, sem \texttt{RetryInfo} & divide o lote \\
    Sobrecarga (\dec{6}) & \texttt{RESOURCE\_EXHAUSTED} + \texttt{RetryInfo} & tenta depois do atraso indicado \\
    Falha transitória & \texttt{UNAVAILABLE}, \texttt{INTERNAL} & tenta mais tarde, com \emph{backoff} \\
    Chave inválida & \texttt{UNAUTHENTICATED} & pausa até a configuração ser corrigida \\
    \bottomrule
  \end{tabularx}
\end{table}

\paragraph{Teoria.}
\emph{Manage Interfaces} (interoperabilidade): o contrato passa a ser um artefacto próprio, com sintaxe e
semântica explícitas. \emph{Reduce Coupling} e \emph{Encapsulate} (modificabilidade): as duas implementações
deixam de depender uma da outra e passam a depender de uma interface versionada, como os repositórios
dependem da \emph{tag} \texttt{v1} do pipeline (secção~\ref{sec:pipeline}).

\paragraph{Alternativas.}
A decisão original de \dec{11} era manter HTTP/JSON, acrescentar um campo \texttt{code} estável e
publicar uma especificação OpenAPI com exemplos testados dos dois lados. Resolvia M2, mas a
especificação ficava ao lado do código: nada obrigava as implementações a segui-la, a não ser os testes.
Com Protocol Buffers é o compilador que obriga. Inicialmente o gRPC tinha sido rejeitado como ``mudança
de tecnologia sem ASR''; essa rejeição foi revista porque o ASR existe (\cen{MO2}), e o desempenho
\emph{não} foi argumento: a ingestão nominal são 33 pedidos por segundo (secção~\ref{sec:estimativas}).
Também se rejeitou pôr o \texttt{.proto} dentro do central (o remote passaria a depender das versões do
central), copiá-lo para cada repositório (duas cópias divergem, que é o próprio problema M2) e manter
HTTP e gRPC em paralelo (duplicaria a autenticação, a validação e o controlo de admissão sem haver
agentes antigos em produção).

\paragraph{Trade-offs.}
Há mais um repositório e um passo de geração de código, mitigado por o código gerado ficar versionado no
\texttt{monitor-api}. O mínimo da linguagem subiu para Go~1.26, exigido pela biblioteca gRPC. Um
balanceador entre remote e central tem de suportar HTTP/2 (\dec{8}), e os testes de carga usam o
cliente gRPC do k6 (\dec{10}). No gRPC a mensagem é lida antes de o pedido chegar ao código da aplicação,
o que permitiria a um cliente sem chave obrigar o central a ler até 1~MiB; para manter a proteção que o
HTTP já tinha, um pedido sem chave plausível é rejeitado ao nível do transporte, antes de a mensagem ser
lida, e só depois a chave é verificada na base de dados.

\paragraph{Verificação.}
\begin{itemize}
  \item \cen{MO2}: o remote decide apenas pelo código e pelo \texttt{RejectionReason}. Um teste envia uma
        rejeição com uma mensagem enganadora e confirma que o motivo guardado na quarentena vem só do
        código. Uma alteração incompatível ao \texttt{.proto} faz falhar o CI do \texttt{monitor-api}
        (\texttt{buf breaking}), que passa a ser a medida executável de \cen{MO2}.
  \item \cen{AV1}, numa versão reduzida: com o central parado a meio de um teste ponta-a-ponta (central
        real com PostgreSQL, remote e SDK), as aplicações continuaram a registar; quando o central voltou,
        chegaram os 40 erros gerados, com 0 duplicados e 0 registos pendentes ou em quarentena.
  \item Métricas de código no remote (iteração~2 contra iteração~1): \emph{issues} de lint de 29 para
        24, funções acima do limite ciclomático de 3 para 1 e acima do limite cognitivo de 6 para 5. A
        classificação das respostas, que antes interpretava códigos HTTP e textos, passou a uma
        correspondência direta de códigos.
\end{itemize}

\porescrever{\dec{12}--\dec{15}.}

\subsection{Disponibilidade e monitorabilidade}
\porescrever{\dec{16}--\dec{17}.}

\subsection{Segurança}
\porescrever{\dec{18}.}

\subsection{Processo e qualidade}
\porescrever{\dec{1}--\dec{3}.}

\section{Avaliação}
\label{sec:avaliacao}

\subsection{Método de medição}
\label{sec:metodo}

As medições de carga correm num ambiente próprio, descrito em código no ficheiro
\texttt{loadtest/compose.yaml}: o
central e o PostgreSQL~16 com \textbf{2 vCPU e 4~GB cada}, e o Prometheus a recolher as métricas do central a
cada 5~s (\dec{4}). A máquina é um AMD Ryzen~7 3700X com 16 \emph{threads} e 31~GB de memória. Cada medição
recria o ambiente do zero, com a base de dados vazia e o central acabado de arrancar, para nenhuma herdar
dados nem \emph{caches} da anterior. As medições não correm no CI, porque os \emph{runners} partilhados dão
números com demasiado ruído.

\paragraph{Agentes simulados.}
O k6 simula 1\,000 \emph{hosts}, um por utilizador virtual, cada um com a sua chave e a sua ligação gRPC. Cada
um comporta-se como o \texttt{monitor-remote}: envia estatísticas a cada 30~s e erros a cada 10~s, em lotes
de até 100 registos, tenta de novo com \emph{backoff} nas falhas transitórias e desiste nas definitivas. Os
erros seguem o perfil da Tabela~\ref{tab:perfil-logs}, atribuído de forma determinística pelo índice do
\emph{host}. Simular os agentes em vez de correr 1\,000 processos \texttt{monitor-remote} é o que torna a
medição possível numa só máquina; a simulação usa o mesmo contrato (\texttt{.proto} do \texttt{monitor-api})
e as mesmas regras de envio.

\paragraph{Três fontes, cada uma para o que mede bem.}
O k6 dá a visão do agente: latência de cada pedido, atraso acrescentado pelo central (definido na
secção~\ref{sec:requisitos}) e lotes que o agente teria posto de parte. O Prometheus dá a do servidor:
percentis da ingestão e do \emph{dashboard}, CPU, memória e ocupação do \emph{pool} de ligações. A base de dados
dá a integridade: no fim de cada medição compara-se o número de registos confirmados ao agente com o número
gravado, e procuram-se duplicados. As contagens vêm do k6 e da base de dados, e não dos contadores do
Prometheus, que perdem o que acontece antes da primeira recolha.

\paragraph{Medições.}
\texttt{nominal} (\cen{SC1}) corre 10~minutos com o perfil de logs; \texttt{nominal-x2} e
\texttt{nominal-x5} repetem-no com a taxa de logs multiplicada por 2 e por 5, porque o perfil é uma hipótese
(análise de sensibilidade). \texttt{noisy-neighbour} (\cen{SC5}) põe 1\,\% dos \emph{hosts} a 200 logs/s durante
2~minutos, a meio de uma medição nominal. \texttt{recovery-storm} (\cen{SC2}) dá a cada \emph{host} 1~h de
\emph{backlog} (cerca de 441\,000 registos no total) e põe-nos todos a enviar ao mesmo tempo, enquanto o
\emph{dashboard} é consultado a cada 2~s. \texttt{capacidade-1}, \texttt{-2} e \texttt{-3} (\cen{SC3}) põem a mesma
carga acima da capacidade com 1, 2 e 3 instâncias atrás do nginx. \texttt{dashboard-30d} (\cen{SC4}) gera
primeiro 31 dias de histórico e corre a carga nominal por cima, com as consultas do \emph{dashboard}, em duas
fases: sem e com a retenção (secção~\ref{sec:sc4}). Desde \dec{9}, a retenção está ligada por omissão, por isso
as medições finais correm com ela ativa, ao contrário das iterações 3 a 5; sem dados expirados, cada execução
não apaga nada e custa poucos milissegundos.

\paragraph{Arquitetura e níveis de teste.}
A arquitetura define as unidades e as interfaces que cada nível de teste exercita
(secção~\ref{sec:documentacao}). A Tabela~\ref{tab:niveis-teste} diz em que nível se verifica cada grupo de
cenários.

\begin{table}[htbp]
  \centering\small
  \caption{Cenários de qualidade por nível de teste}
  \label{tab:niveis-teste}
  \begin{tabularx}{\linewidth}{L{0.2\linewidth}L{0.3\linewidth}Y}
    \toprule
    Nível & O que exercita & Cenários \\
    \midrule
    Unitário & uma unidade da vista de módulos, com \emph{fakes} das interfaces (\dec{14}) & \cen{TE2}, \cen{TE3}, \cen{MO1}, \cen{MO2}, \cen{AV3}, \cen{AV4}, \cen{MN2} \\
    Integração & as interfaces entre elementos: os repositórios num PostgreSQL real, o contrato (\texttt{buf}), o SDK contra o esquema do remote & \cen{AV2}, \cen{SC3} (liderança), \cen{SE1}, \cen{SE3}, \cen{MO2} \\
    Aceitação de QA: carga & o sistema inteiro sob carga, com os limites de recursos (k6) & \cen{SC1}--\cen{SC5}, \cen{IO1} \\
    Aceitação de QA: ponta a ponta & remote, SDK, nginx com TLS e central reais em Compose & \cen{AV1}, \cen{MN1}, \cen{SE1} \\
    Inspeção & \emph{checklist}, regras estáticas, experiência controlada & \cen{SE2} (ASVS), \cen{MO3}, camadas (\texttt{depguard}) \\
    \bottomrule
  \end{tabularx}
\end{table}

\paragraph{Limitações.}
O gerador de carga corre na mesma máquina que o sistema medido. Os limites de CPU do central e da base de
dados impedem que o k6 lhes tire recursos, mas não isolam totalmente a memória e o disco partilhados. Os
números absolutos dependem desta máquina; o que se compara são as iterações entre si, medidas da mesma
forma.

\subsection{Escalabilidade: antes e depois}
\label{sec:aval-escala}

As tabelas desta secção comparam duas iterações, medidas da mesma forma (secção~\ref{sec:metodo}): a
iteração~3, com o contrato gRPC e as métricas mas antes de qualquer decisão de escalabilidade (a
\emph{baseline} de carga), e a iteração~7, a versão final, com todas as decisões e a retenção ligada por
omissão. O texto cita também a iteração~4, a primeira medição depois de \dec{5} (escritas em lote) e \dec{6}
(controlo de admissão), para separar o efeito destas duas decisões do das seguintes.

\gerado{carga-sc1.tex}

\paragraph{Regime normal (\cen{SC1}).}
Já na \emph{baseline}, \cen{SC1} cumpria-se com uma folga muito grande (Tabela~\ref{tab:carga-sc1}): p95 da
ingestão de 6~ms contra um limite de 300~ms, e o central a usar 4\,\% dos 2 vCPU. O débito medido, 122 registos
por segundo, coincide com a estimativa feita antes de qualquer medição (cerca de 90 logs e 33 estatísticas por
segundo, secção~\ref{sec:estimativas}): a medição confirmou a estimativa e a conclusão que se tirou dela, de que
a ingestão nominal não era o problema de escala. Na versão final, o p95 é de 5~ms e o CPU de 3\,\%.

\gerado{carga-sensibilidade.tex}

\paragraph{Sensibilidade ao perfil de logs.}
Como o perfil de logs é uma hipótese, cada medição nominal repetiu-se com o dobro e o quíntuplo da taxa de logs
(Tabela~\ref{tab:carga-sensibilidade}). Na \emph{baseline}, com cinco vezes mais logs, o p95 subia para 10~ms e o
CPU para 8\,\%; na versão final ficam em 5~ms e 4\,\%, praticamente iguais ao regime normal, porque o custo de
um lote deixou de crescer com o número de registos. \cen{SC1} cumpre-se em todos os casos: a conclusão não
depende de a hipótese estar certa.

\gerado{carga-sc5.tex}

\paragraph{Vizinho ruidoso (\cen{SC5}).}
Enquanto 1\,\% dos \emph{hosts} envia 200 logs por segundo, o atraso p95 dos restantes foi de 6~ms na
\emph{baseline} e de 5~ms depois (Tabela~\ref{tab:carga-sc5}). O risco que motivava a parte de justiça entre
\emph{hosts} de \dec{6} (lacuna S10) não se materializou a esta escala. É o \emph{generate and test} a funcionar: a
medição refutou uma hipótese da análise, e essa parte de \dec{6} ficou desenhada mas não implementada
(secção~\ref{sec:d6}). Os próprios \emph{hosts} em rajada também ganharam: o seu atraso p95 desceu de 43 para 11~ms.

\gerado{carga-sc2.tex}

\paragraph{Recuperação (\cen{SC2}).}
Era o único cenário que falhava, e falhava onde a análise previa (Tabela~\ref{tab:carga-sc2}). Na
\emph{baseline}, os 1\,000 \emph{hosts} entregavam as 441\,000 linhas do \emph{backlog} em 1,4~minutos, sem perdas
nem duplicados, mas o \emph{dashboard} ficava inutilizável, com p95 de pelo menos 10~s (o limite superior do
histograma). O CPU do central nunca passou de 30\,\%; o recurso esgotado era o \emph{pool} de ligações à base de
dados (25 de 25), com mais de 11\,000 esperas por uma ligação livre, e 9,5\,\% das tentativas de envio
acabavam em \emph{timeout}.

Depois de \dec{5} e \dec{6} (iteração~4), o \emph{dashboard} passou a ter p95 de 262~ms durante a recuperação, e
\cen{SC2} passou a cumprir-se; na versão final, com o \emph{bulkhead} de \dec{9}, é de 145~ms. A ingestão nunca ocupa mais de 15 ligações (o limite de \dec{6}), e já não há esperas nem
\emph{timeouts}. Quase um quarto das tentativas (24,8\,\% na versão final) é rejeitado por sobrecarga, mas são
rejeições intencionais: o agente volta a enviar depois do atraso indicado, e as 441\,000 linhas chegam todas, sem
duplicados. O resultado menos óbvio é que a recuperação ficou \emph{mais rápida}, 42~s em vez de 1,4~minutos,
apesar das rejeições. Com escritas em lote cada pedido ocupa uma ligação durante muito menos tempo, e limitar quantos
pedidos competem pela base de dados ao mesmo tempo elimina a contenção que, antes, fazia cada um demorar
segundos. O \emph{trade-off} previsto em \dec{6} (recuperação mais lenta em troca de estabilidade) não se
verificou: as duas decisões juntas melhoraram ambos os critérios.

\paragraph{Regressão (iterações~5 e~7).}
Depois de \dec{8}, as medições de \cen{SC1}, \cen{SC2}, \cen{SC5} e da sensibilidade repetiram-se com uma
instância. Todos os cenários continuaram a cumprir-se, com valores dentro da variação entre execuções (o p95 do
\emph{dashboard} durante a recuperação foi de 380~ms, contra 262~ms na iteração~4; o limite é 1~s). A medição
final (iteração~7) repetiu todos os cenários de carga sobre a versão entregue, com a retenção ligada, e nenhum
regrediu. A retenção corre também durante \cen{SC2}, mas sem dados expirados cada execução custa
milissegundos (p95 de 9,5~ms).

\subsection{Escala horizontal (\cen{SC3})}
\label{sec:aval-sc3}

\gerado{carga-sc3.tex}

\cen{SC3} mede o débito com 1, 2 e 3 instâncias atrás do nginx, com a mesma carga: os 1\,000 \emph{hosts} enviam
lotes de 100 erros sem pausa (Tabela~\ref{tab:carga-sc3}). A parte de coordenação cumpre-se: em todas as
medições, cada \emph{job} teve um único líder e correu numa só instância, e no ensaio da secção~\ref{sec:d8} a
troca de líder demorou 12~s, contra um limite de 30~s. A parte do débito \textbf{não se cumpre}: duas e três
instâncias confirmaram 1,01 e 1,04 vezes os registos de uma na medição final (1,01 e 1,03 na iteração~5),
contra os 1,8 e 2,5 vezes pedidos.

A tabela mostra porquê. Cada instância do central usa em média 11\,\% ou menos do seu CPU; o que sobe com as
instâncias é a latência de cada pedido (p95 de 100, 199 e 312~ms), quase na proporção da concorrência admitida
(15, 30 e 45 pedidos). Pela lei de Little, débito igual com o triplo da concorrência e o triplo da latência quer dizer que os
pedidos passam a esperar num recurso partilhado: o PostgreSQL. Uma amostragem dos processos ativos do
PostgreSQL durante a carga confirma-o: 56\,\% das amostras estavam a executar, 26\,\% à espera do registo de
transações (WAL: inserção, escrita e sincronização) e 11\,\% em contenção nas mesmas páginas de dados. O limite
é o caminho de escrita da base de dados, e não o central.

A premissa de \cen{SC3}, uma carga acima da capacidade de uma instância, deixou de existir com \dec{5}. As
escritas em lote tornaram o central tão leve que o estrangulamento passou para a base de dados. Uma instância
confirmou cerca de 22\,000 registos por segundo no minuto de pico, cerca de 180 vezes os 122 registos por segundo
do regime normal. Este fator é um limite superior: foi medido com lotes cheios, o caso mais eficiente, e os
\emph{hosts} reais enviam lotes pequenos. Mesmo assim, a folga cobre com margem os 10\,000 \emph{hosts} do
\emph{backlog} do projeto (cerca de 1\,200 registos por segundo).

Por isso, a equipa decidiu não escalar a base de dados. Réplicas de leitura não aumentam o débito de escrita;
\emph{sharding} seria uma complexidade que nenhum cenário pede; e desligar a sincronização do WAL poria em risco
\cen{AV1}, porque o agente apaga do \emph{buffer} o que o central confirma. Escalar a base de dados fica como
evolução E9 (secção~\ref{sec:evolucao}), com um gatilho medido: débito sustentado acima de metade da capacidade aqui medida. \dec{8} mantém
o seu valor pela disponibilidade: o central deixa de ser um ponto único de falha, e um \emph{job} sobrevive à
morte da instância que o corria. É o \emph{generate and test} outra vez: a medição mostrou que a premissa de um
cenário deixou de ser verdadeira, e o relatório regista-o em vez de ajustar o alvo.

Na iteração~5, o \emph{dashboard} também piorava sob esta carga extrema (p95 de 4,9~s com uma instância e de
9,4~s com três): o limite de admissão de \dec{6} é por instância, e três instâncias metem três vezes mais
pedidos em simultâneo na base de dados. Na versão final fica entre 84 e 174~ms: em cada instância, a ingestão
(15 ligações) e o \emph{dashboard} (5, o \emph{bulkhead} de \dec{9}) têm compartimentos separados dentro do
\emph{pool} de 25, e as consultas, agora de custo limitado, libertam a ligação depressa. A regra de configuração mantém-se: com N instâncias, \texttt{INGEST\_MAX\_INFLIGHT} deve ser o orçamento
total da ingestão a dividir por N (ADR-0007).

\subsection{Trinta dias de dados (\cen{SC4})}
\label{sec:sc4}

\gerado{carga-sc4.tex}

A Tabela~\ref{tab:carga-sc4} compara a medição de \cen{SC4} antes de \dec{9} com a medição final, cada uma em
duas fases de 10~minutos sobre 31 dias de histórico (89~M estatísticas, 239~M erros, 143 a 150~GB). Antes, a
primeira fase corria sem retenção e a segunda com a retenção a 30 dias a cada minuto; na versão final a
retenção está ligada por omissão, e as duas fases são duas corridas com ela. \textbf{Antes, \cen{SC4} falhava nas duas partes.} As sete consultas do
\emph{dashboard} tinham p95 entre 29 e 74~s, e a ingestão, mesmo sem retenção, confirmava só cerca de 11\,000
dos 73\,000 registos esperados em 10~minutos, com 58\,\% dos pedidos rejeitados por sobrecarga e um atraso
acrescentado pelo central de 636~s no p95. Nenhum registo se perdeu: os que não entraram ficaram nos agentes,
por entregar, como \dec{19} e o ADR-0020 preveem. A causa foi isolada medindo cada consulta sozinha
(secção~\ref{sec:d9}): quatro consultas sem limite temporal ocupavam todo o \emph{pool} de ligações.

\textbf{Depois, \cen{SC4} cumpre-se.} Na medição final, o p95 das sete consultas fica entre 15 e 210~ms (o pior
é o das métricas de um mês de um \emph{host}; na iteração~6, a primeira medição depois de \dec{9}, era de
287~ms), e a ingestão sobre 30 dias de dados fica ao nível da medição nominal com a base vazia: p95 de 5~ms,
nenhum pedido rejeitado, todos os registos confirmados, sem duplicados. O CPU do PostgreSQL desceu de 67\,\%
para 7\,\% em média. Isoladas, as quatro consultas passaram de 55~s, 17~s, 13~s e
5,3~s para 19~ms, 0,6~ms, 2~ms e 265~ms (estas últimas com a \emph{cache} fria).

\paragraph{A retenção e a decisão de não particionar.}
Antes, cada execução da retenção demorava até 27~s, perto do \emph{timeout}, não pelo seu custo, mas porque
esperava pelas ligações que o \emph{dashboard} ocupava. Depois, cada execução demora 0,2 a 0,3~s em média
(0,48~s no p95, nas três corridas com retenção) e o número de registos expirados desce cerca de 733 mil por
hora: a retenção remove 2,6 vezes o que entra, e o atraso acumulado esvazia-se. O p95 da ingestão não subiu com
a retenção (na iteração~6, 7,4~ms sem ela e 4,9~ms com ela; a fase com retenção corre depois e encontra a
\emph{cache} quente, e a diferença é da ordem do ruído). Nenhuma das três condições escritas no ADR-0017 antes da medição se
verifica, por isso \textbf{o particionamento não se fez}. Fica como o passo seguinte, com um gatilho sobre
métricas que já existem: a duração da retenção (p95 acima de 15~s), o p95 da ingestão nos minutos em que a
retenção corre, e o número de registos com mais de 30 dias. É o resultado mais claro do \emph{generate and
test} deste trabalho: o plano previa particionar, e a medição mostrou que o problema estava noutro sítio.

\subsection{Disponibilidade e integridade dos dados}
\label{sec:av1-av4}

\paragraph{\cen{AV2}: \emph{host} sem resposta, com 1.000 \emph{hosts}.}
Um teste de integração, num PostgreSQL real, regista 1.000 \emph{hosts} com \emph{endpoints} num servidor
local; metade nunca responde, e as \emph{probes} esgotam os 5~s de \emph{timeout}. O servidor regista quando cada
\emph{host} recebe cada \emph{probe}. Antes de \dec{7}, o teste corre um ciclo; depois, corre o \emph{job} durante
3~minutos, como em produção. A Tabela~\ref{tab:av2} compara os dois (dados em
\texttt{docs/metrics/iter-6/av2-antes.txt} e \texttt{av2-depois.txt}).

\begin{table}[htbp]
  \centering\small
  \caption{\cen{AV2}: 1.000 \emph{hosts}, metade sem resposta, antes e depois de \dec{7}}
  \label{tab:av2}
  \begin{tabular}{lrrl}
    \toprule
    Medida & Antes & Depois & Alvo \\
    \midrule
    Intervalo entre \emph{probes} do mesmo \emph{host} (máximo) & 350,7~s & 35,4~s & $\le$ 45~s \Sim \\
    Pior atraso até ao alerta, além do limiar & 6~min~6~s & 50,4~s & $\le$ 60~s \Sim \\
    Tempo de cada \emph{host} em ``desconhecido'' por ciclo & 4~min~21~s & 0 & 0 \Sim \\
    Iteração mais longa do \emph{job} & 350,7~s & 20,2~s & $<$ 60~s \Sim \\
    \emph{Probes} por segundo (máximo em simultâneo) & --- & 32,4 (128) & --- \\
    \bottomrule
  \end{tabular}
\end{table}

A mediana do intervalo é 30,3~s e o máximo fica abaixo dos 38~s calculados. A iteração mais longa é a primeira,
com os 1.000 \emph{hosts} devidos ao mesmo tempo; nas seguintes, o \emph{jitter} já os espalhou. O pior atraso
soma ao intervalo máximo uma avaliação de alertas (15~s), sem contar com a fase entre a falha e a próxima
avaliação, que só o reduz. \cen{AV2} cumpre-se.

\paragraph{\cen{AV4}: rajada com o central em baixo.}
Um teste de integração no remote usa um \emph{buffer} SQLite real com um limite de 1~MiB (o rígido a 4~MiB) e
simula 12 ciclos de rajada com o central em baixo: em cada ciclo, uma aplicação escreve 1.000 INFO, 50 WARN e
20 ERROR, e o \emph{worker} de \dec{19} corre uma vez. Cada ciclo escreve metade do limite de uma só vez, muito
mais do que uma rajada real escreve em 10~s, para forçar descartes em quase todos os ciclos.
A Tabela~\ref{tab:av4} resume o resultado (dados em \texttt{docs/metrics/iter-6/av4.txt}).

\begin{table}[htbp]
  \centering
  \small
  \caption{\cen{AV4}: \emph{buffer} de 1~MiB, 12 ciclos de rajada com o central em baixo}
  \label{tab:av4}
  \begin{tabular}{lrl}
    \toprule
    Medida & Valor & Alvo \\
    \midrule
    Logs escritos (INFO / WARN / ERROR) & 12.000 / 600 / 240 & --- \\
    Tamanho depois de cada ciclo (máximo) & 0,88~MiB & $\le$ 1~MiB \Sim \\
    ERROR perdidos & 0 de 240 & 0 \Sim \\
    WARN perdidos & 0 de 600 & só depois dos INFO \Sim \\
    INFO descartados, contados e persistidos & 11.500 & todos contados \Sim \\
    INFO que restam & os 500 mais recentes & mais antigos primeiro \Sim \\
    \bottomrule
  \end{tabular}
\end{table}

O tamanho antes de cada verificação chegou a 1,36~MiB: é o excesso entre duas verificações, aqui exagerado de
propósito (meio limite por ciclo). Com o limite por omissão de 512~MiB e uma rajada real, o excesso estimado é
inferior a 1\,\%. Os contadores sobrevivem a reabrir o ficheiro, como num reinício do agente. \cen{AV4} cumpre-se.

\paragraph{\cen{AV3}: retoma sem reinício.}
Um teste no remote, com o relógio injetado e um central falso em memória, recusa a chave com
\texttt{PERMISSION\_DENIED}. A entrega pausa e, durante 4~min~59~s, nenhuma das duas categorias faz pedidos. O
operador corrige a causa; aos 5~minutos, a \emph{probe} envia um só registo, que leva o estado do \emph{buffer}
com a pausa assinalada, o central aceita-o, e a entrega retoma no mesmo ciclo e entrega tudo, sem reiniciar o
agente. Outros testes fixam que uma \emph{probe} falhada marca a seguinte para 5~minutos depois, que só há uma
\emph{probe} de cada vez, e que o \texttt{SIGHUP} troca a chave e faz a \emph{probe} de imediato. No processo
real, dois \texttt{SIGHUP} seguidos (chave nova e depois um \texttt{.env} sem chave) deixam o agente vivo,
com a chave nova no primeiro caso e a anterior no segundo (dados em \texttt{docs/metrics/iter-6/av3.txt}).
\cen{AV3} cumpre-se.

\paragraph{\cen{AV1}: central em baixo, ponta a ponta.}
O teste ponta a ponta de \dec{15} mata o central (\texttt{docker compose kill}, sem paragem graciosa), com o
balanceador de pé, e a aplicação escreve 500 logs com o SDK durante a falha. O remote continua a recolher
estatísticas e a tentar entregar, recebendo \texttt{UNAVAILABLE}. Quando o central volta, o condutor espera
que o \emph{buffer} esvazie, para o remote para ter um estado estável e compara os dois lados registo a
registo, pelo identificador de origem: tudo o que o remote marcou como entregue tem de estar no central, uma só
vez, e nada no central pode faltar no \emph{buffer}. A comparação abrange o \emph{buffer} inteiro, incluindo os
500 logs do caminho feliz entregues antes da falha. A Tabela~\ref{tab:av1} resume as execuções (dados em
\texttt{docs/metrics/iter-6/e2e/}).

\begin{table}[htbp]
  \centering\small
  \caption{\cen{AV1}: central morto, 500 logs gerados durante a falha}
  \label{tab:av1}
  \begin{tabular}{lrrl}
    \toprule
    Medida & 60~s em baixo & 1~h em baixo & Alvo \\
    \midrule
    Duração real da falha & 63~s & 3.603~s & --- \\
    Pendentes quando o central volta (logs / estatísticas) & 500 / 10 & 500 / 599 & --- \\
    Logs entregues e no central & 1.000 / 1.000 & 1.000 / 1.000 & --- \\
    Estatísticas entregues e no central & 18 / 18 & 634 / 634 & --- \\
    Perdidos & 0 & 0 & 0 \Sim \\
    Duplicados & 0 & 0 & 0 \Sim \\
    No central sem origem no \emph{buffer} & 0 & 0 & 0 \Sim \\
    Pendentes ou em quarentena no fim & 0 & 0 & 0 \Sim \\
    Tempo até entregar tudo, depois de o central voltar & 36,7~s & 202,0~s & --- \\
    \bottomrule
  \end{tabular}
\end{table}

O tempo de recuperação é o do \emph{backoff} do remote: depois de 60~s de falhas seguidas, a espera entre
tentativas já ia em cerca de 40~s (32 a 48~s, com a variação aleatória); o limite do \emph{backoff} é 5~minutos, que é por isso o
atraso máximo depois de uma falha longa. Com 1~h em baixo, a entrega completa demorou 202~s: até 5~minutos de
espera pela tentativa seguinte, e depois 500 logs e 599 estatísticas em lotes. As duas execuções foram feitas com
o código ainda por \emph{commit}; o \emph{workflow} de ponta a ponta (\texttt{e2e.yml}) repete-as a partir do
\texttt{main}. Os duplicados são impossíveis no central por construção (recibos por
host e identificador de origem); o que o teste prova é que o remote não cria identificadores novos ao repetir
um envio e que não marca como entregue nada que o central não tenha gravado. \cen{AV1} cumpre-se.

\subsection{Manutenção, testabilidade e interoperabilidade: MO1--MO3, TE2, TE3, IO1 e métricas de código}
\label{sec:mo1}

\paragraph{\cen{MO1}: uma métrica de alerta nova.}
Com o mecanismo de \dec{13} no lugar, implementámos uma métrica real que o projeto precisava: o número de
registos em quarentena reportado por cada agente (\texttt{buffer\_quarantined}, \dec{17}). A Tabela~\ref{tab:mo1}
compara o que a alteração tocou com a linha de base contada antes de \dec{13}: 7 sítios em 5 ficheiros (três
Go, um modelo HTML e a migração), a que se somam os dois ficheiros Go que qualquer migração obriga a mudar
(dados em \texttt{docs/metrics/iter-6/mo1.md}).

\begin{table}[htbp]
  \centering\small
  \caption{\cen{MO1}: o que uma métrica de alerta nova exige, antes e depois de \dec{13}}
  \label{tab:mo1}
  \begin{tabular}{lccl}
    \toprule
    Medida & Antes & Depois & Alvo \\
    \midrule
    Ficheiros Go com código da métrica & 3 & 1 & --- \\
    Ficheiros Go que registam a migração & 2 & 2 & --- \\
    Ficheiros Go de produção, no total & 5 & 3 & $\le$ 3 \Sim \\
    Migrações & 1 & 1 & 1 \Sim \\
    Modelos HTML & 1 & 0 & --- \\
    Sítios a mudar fora da métrica e da migração & 6 & 0 & 0 \Sim \\
    Camadas tocadas & 3 & 1 & --- \\
    \bottomrule
  \end{tabular}
\end{table}

Dos três ficheiros Go, só um é da métrica: a estratégia, com 36 linhas. Os outros dois registam a migração e
mudam em qualquer migração, de métricas ou não. Um executor que descobrisse as migrações pelo nome do ficheiro
tirá-los-ia, mas a ordem explícita das migrações é uma garantia que preferimos manter. Os testes tocados foram
dois: o do registo e o de caracterização, que passou a fixar também o valor da métrica nova numa base de
dados real. O teste que compara o registo com o \texttt{CHECK} da base de dados confirmou que a migração
acompanha a métrica. A alteração foi feita com um assistente de programação, em pouco mais de um minuto de
relógio; esse tempo não representa o de uma pessoa sem assistência, mas a dimensão da alteração (um ficheiro
novo a seguir um modelo que já existe) deixa folga larga para o alvo de uma pessoa-dia. \cen{MO1} cumpre-se.

\paragraph{\cen{MO2}: uma definição por repositório.}
A Tabela~\ref{tab:mo2} conta, antes e depois de \dec{12}, quantas vezes a lista dos níveis de log está escrita em
Go, por repositório (dados em \texttt{docs/metrics/iter-6/te2-te3-mo2.md}). As cópias que ficam no SQL e nas
traduções do \emph{protobuf} são cobertas por seis testes de consistência. Para confirmar que apanham uma
divergência, retirámos um nível do \texttt{CHECK} do central e, noutra vez, da lista do SDK: em ambos os casos o
teste respetivo falhou, e a alteração foi revertida. No remote, as duas validações do endereço do central
passaram a uma. \cen{MO2} cumpre-se.

\begin{table}[htbp]
  \centering\small
  \caption{\cen{MO2}: lista dos níveis de log escrita em Go, por repositório}
  \label{tab:mo2}
  \begin{tabular}{lccl}
    \toprule
    Repositório & Antes & Depois & Verificada contra \\
    \midrule
    central & 5 & 1 & \texttt{CHECK} do SQL, \emph{protobuf}, SQL das métricas \\
    remote & 1 & 1 & \texttt{CHECK} do \emph{buffer}, \emph{protobuf} \\
    SDK & 1 (privada) & 1 (exportada) & \texttt{CHECK} do \emph{buffer} \\
    hello-world & 1 & 0 & usa a do SDK \\
    \bottomrule
  \end{tabular}
\end{table}

\paragraph{\cen{MO3}: uma mudança de apresentação.}
Medimo-lo com uma experiência controlada, numa cópia do central descartada no fim (dados e \emph{diff} em
\texttt{docs/metrics/iter-6/mo3.md}). Fizemos duas alterações de apresentação típicas, escolhidas antes de olhar
para o código que as serve. No \emph{dashboard}, os incidentes abertos passaram a ser o primeiro cartão, e o
cartão de erros passou a mostrar quantos dos erros de hoje são críticos. Na página de erros, o nível passou a
\emph{badge} com uma cor por severidade, e o ambiente ganhou uma coluna própria. As duas tocaram quatro
ficheiros, todos em \texttt{frontend/} (um modelo HTML por página, o JavaScript do \emph{dashboard} e a folha de
estilos), e nenhum ficheiro Go: \textbf{0 alterações em \texttt{services} e \texttt{repositories}}. Os testes
dos \emph{handlers}, que desenham os modelos reais, passaram sem alterações. \cen{MO3} cumpre-se. Não usámos
o histórico do git como medida porque os \emph{commits} são por iteração e juntam alterações de várias
naturezas. O limite da medida é o do próprio cenário: uma página que precise de um dado novo já não é só
apresentação, e toca nas camadas de baixo por desenho.

\paragraph{\cen{IO1}: um agente noutra linguagem.}
O gerador de carga do D10 é esse agente: está escrito em JavaScript, para o k6, e o único artefacto do sistema
que usa é o \texttt{.proto} do \texttt{monitor-api}, montado só de leitura (dados em
\texttt{docs/metrics/iter-6/io1.md}). Tudo o que precisa de saber sobre o comportamento está nos comentários do
contrato: a autenticação, a deduplicação por \texttt{source\_id}, o que fazer com cada código gRPC e com o
\texttt{RetryInfo}. A lógica de agente ocupa 101 linhas. Com as mesmas garantias do remote, confirmou 73\,229
registos na medição de \cen{SC4}, todos gravados, sem duplicados nem lotes perdidos, e fez o mesmo nos
441\,000 registos de \cen{SC2}. \cen{IO1} cumpre-se, com dois limites: a durabilidade local (o \emph{buffer}) é
responsabilidade de cada agente e não do contrato, e o agente foi escrito pela mesma equipa que escreveu o
contrato.

\paragraph{\cen{TE1}: integração contínua.}
Nos últimos 20 \emph{runs} com sucesso de cada repositório (lidos no GitHub a 21/11), o CI demorou em mediana
entre 89~s (\texttt{monitor-api}) e 141~s (\texttt{monitor-central}), e no máximo 401~s (\texttt{monitor-hello-world}),
abaixo dos 10~minutos (alvo). Uma falha bloqueia o \emph{merge}: o \texttt{main} de cada repositório está
protegido, também para os administradores, e aconteceu durante o trabalho. A Higiene do PR de \dec{5} e
\dec{6} falhou por falta de um \texttt{go mod tidy}, e o \emph{quality gate} do SonarQube recusou os PR do gRPC,
de \dec{4} e de \dec{5}--\dec{6} por cobertura do código novo abaixo de 80\,\%; nenhum entrou antes de ser
corrigido. \cen{TE1} cumpre-se.

\paragraph{\cen{TE2} e \cen{TE3}: \emph{handlers} e tempo.}
Com as interfaces de \dec{14}, a cobertura dos \emph{handlers} passou de 9,5\,\% para 95,6\,\% (alvo
$\ge$ 60\,\%). Os testes cobrem todos os \emph{handlers}, desenham as páginas reais e correm em cerca de 0,2~s,
sem base de dados. Nos testes dos quatro repositórios, os \texttt{time.Sleep} passaram de 1 para 0 (alvo 0). O
resumo diário é testado em seis casos com o relógio e o fuso controlados: antes, às e depois das 9~h, no fim do
mês, no dia em que a hora muda, e noutro fuso. \cen{TE2} e \cen{TE3} cumprem-se.

\gerado{metricas-codigo-antes-depois.tex}

\subsection{Monitorabilidade}
\label{sec:mn}

\paragraph{\cen{MN1}: quarentena visível no central.}
No remote, um teste põe um registo em quarentena com 150 estatísticas pendentes e verifica que os dois lotes do
ciclo levam a mesma fotografia, com a quarentena contada. No central, um teste de integração em PostgreSQL envia
um lote com três registos em quarentena e lê-os na API do \emph{dashboard}; verifica também que uma fotografia
mais antiga não substitui a mais recente e que uma inválida é ignorada sem perder as estatísticas. A quarentena
fica visível no lote seguinte.

O teste ponta a ponta de \dec{15} repete-o com os componentes reais: corrompe um log pendente no \emph{buffer}
(metadados que deixam de ser JSON), espera que o remote o ponha em quarentena e mede quanto tempo demora a
aparecer no central. Na primeira execução falhou, com 10,0~s para ciclos de 5~s: um ciclo de entrega que acorda
antes da recolha não tem estatísticas para levar a fotografia, e o central ficava dois ciclos atrás
(secção~\ref{sec:d15}). Depois da correção, três execuções mediram 5,85~s, 0,16~s e 5,82~s. O atraso depende da fase
entre a quarentena e a recolha seguinte, e o máximo é um ciclo de recolha, que no teste são 6~s: o intervalo
mais 1~s de amostragem do CPU. O limite do teste é 7~s, com 1~s para a observação pelo condutor; em produção,
o ciclo é de 31~s. \cen{MN1} cumpre-se, nos testes de cada lado e ponta a ponta.

\paragraph{\cen{MN2}: \emph{job} parado.}
Um teste corre um \emph{job} liderado com intervalo de 15~s que bloqueia, com o relógio controlado: aos 30~s o
\texttt{/health/ready} ainda responde 200, um segundo depois responde 503 e nomeia o \emph{job}, e o
\texttt{/health/live} continua a responder 200. Um \emph{job} liderado por outra instância nunca conta. Quando
a iteração termina, a instância volta a estar pronta. A deteção acontece em $\le$ 2 intervalos. \cen{MN2}
cumpre-se (dados em \texttt{docs/metrics/iter-6/mn1-mn2.txt}).

\subsection{Segurança}
\label{sec:aval-se}

\paragraph{\cen{SE1}: chave inválida ou revogada.}
O teste ponta a ponta arranca um segundo remote, com o seu próprio \emph{buffer}, duas vezes: com uma chave que
não existe e com a chave antiga de um host cuja chave foi regenerada (a mesma alteração que a interface faz:
um \emph{digest} novo na linha do host). Em cada caso, a aplicação escreve 10 logs. Nos dois casos, o remote
pausa por \texttt{UNAUTHENTICATED} em 0,7~s (\dec{16}); o central não tem nenhuma linha de logs, estatísticas
ou estado do \emph{buffer} de outro host que não o principal, e o host revogado nunca aparece como visto; no
\emph{buffer} do remote, os 10 logs e a estatística ficam pendentes, sem quarentena: corrigida a chave,
seriam entregues. \cen{SE1} cumpre-se.

\paragraph{\cen{SE2}: XSS e \emph{clickjacking}.}
Um teste verifica que todas as respostas, de página, de API, de erro e de rota inexistente, levam a CSP e os outros
cabeçalhos, e o teste das páginas falha se alguma tiver um \emph{script inline} ou de uma CDN. A prova principal
foi feita num navegador real, com a CSP em \emph{enforce}: as sete páginas da interface carregaram sem nenhum
relatório de violação; um \emph{script} injetado na página, como faria um XSS, não executou; e a página posta num
\texttt{<iframe>} foi bloqueada. Os dois bloqueios chegaram ao \emph{endpoint} de relatórios do central
(dados em \texttt{docs/metrics/iter-6/se2-se3.md}). \cen{SE2} cumpre-se.

\paragraph{\cen{SE3}: \emph{endpoint} de saúde para \emph{loopback} ou metadados.}
A guarda das \emph{probes} é testada com 14 endereços que têm de ser recusados (\emph{loopback}, metadados das
\emph{clouds}, \emph{link-local}, \emph{multicast}, endereços IPv4 escritos como IPv6) e 6 que têm de passar
(redes privadas e públicas). Num teste ponta a ponta, \emph{probes} reais para \texttt{127.0.0.1} e para
\texttt{localhost}, um nome que só depois de resolvido se sabe que é \emph{loopback}, não chegaram ao servidor
local, e os dois \emph{hosts} ficaram em baixo; sem a guarda, o mesmo servidor recebeu a \emph{probe}, o que mostra
que foi a guarda a travá-la. \cen{SE3} cumpre-se.

\paragraph{OWASP ASVS 5.0.0, nível 1.}
Dos 70 requisitos, antes de \dec{18} cumpriam-se 36, 25 não se aplicavam (não há \emph{upload} de ficheiros,
OAuth, \emph{tokens} autocontidos, XML nem WebSockets) e 9 falhavam; depois, cumprem-se 42 e falham 3, com o risco
aceite e um gatilho no ADR-0016 (avaliação completa em \texttt{docs/seguranca/asvs-n1.md}).

\subsection{Resumo dos cenários}
\label{sec:resumo-cenarios}

A Tabela~\ref{tab:resumo-cenarios} junta os 21 cenários do capítulo~\ref{sec:requisitos}, com a medida-alvo, o
valor medido e o veredicto. Os valores de carga são os da medição final (iteração~7); os restantes, das
medições indicadas em cada secção. Cumprem-se 20; \cen{SC3} cumpre-se na coordenação e não no débito, pela razão
da secção~\ref{sec:aval-sc3}.

\begin{table}[htbp]
  \centering\small
  \caption{Resumo dos cenários: alvo, valor medido e veredicto}
  \label{tab:resumo-cenarios}
  \begin{tabularx}{\linewidth}{lYYc}
    \toprule
    Cenário & Medida-alvo & Medido & Cumpre \\
    \midrule
    \cen{SC1} & p95 $<$ 300~ms; 0\,\% de erros; CPU $<$ 70\,\% & 5~ms; 0\,\%; 3\,\% & \Sim \\
    \cen{SC2} & \emph{backlog} em $\le$ 15~min; \emph{dashboard} p95 $<$ 1~s; 0 perdas e duplicados & 42~s; 145~ms; 0 e 0 & \Sim \\
    \cen{SC3} & débito $\ge$ 1,8× e $\ge$ 2,5×; 0 \emph{jobs} duplicados; troca de líder $\le$ 30~s & 1,01× e 1,04×; 0; 12~s & \Parc \\
    \cen{SC4} & \emph{dashboard} p95 $<$ 1~s com 30 dias; ingestão sem degradação com a retenção & 15 a 210~ms; ingestão p95 5~ms, 0 rejeições & \Sim \\
    \cen{SC5} & atraso p95 dos outros \emph{hosts} $<$ 5~s & 5~ms & \Sim \\
    \addlinespace
    \cen{AV1} & 1~h em baixo: 0 perdas, 0 duplicados & 3.603~s: 0 e 0 & \Sim \\
    \cen{AV2} & alerta em $\le$ \emph{timeout} + 60~s com 1\,000 \emph{hosts} & 50,4~s além do limiar & \Sim \\
    \cen{AV3} & retoma em $\le$ 5~min sem reiniciar & aos 5~min, sem reinício & \Sim \\
    \cen{AV4} & \emph{buffer} $\le$ limite; só INFO e WARN perdidos, contados & 0,88 de 1~MiB; 0 ERROR e 0 WARN perdidos; 11.500 INFO contados & \Sim \\
    \addlinespace
    \cen{MO1} & $\le$ 1 pessoa-dia; $\le$ 3 ficheiros Go e 1 migração & 3 ficheiros Go (1 da métrica) e 1 migração & \Sim \\
    \cen{MO2} & 1 sítio por repositório; divergência apanhada por teste & 1 por repositório; 2 divergências, 2 apanhadas & \Sim \\
    \cen{MO3} & 0 alterações em \texttt{services} e \texttt{repositories} & 0 (4 ficheiros, todos em \texttt{frontend/}) & \Sim \\
    \addlinespace
    \cen{TE1} & CI em $\le$ 10~min; uma falha bloqueia o \emph{merge} & mediana 89 a 141~s, máximo 401~s; bloqueou 3 PR & \Sim \\
    \cen{TE2} & cobertura dos \emph{handlers} $\ge$ 60\,\% & 95,6\,\% & \Sim \\
    \cen{TE3} & lógica temporal sem esperas & 0 \texttt{time.Sleep} nos testes & \Sim \\
    \addlinespace
    \cen{MN1} & quarentena visível em $\le$ 1 ciclo & no máximo 5,85~s, com ciclos de 6~s & \Sim \\
    \cen{MN2} & \texttt{/health/ready} falha em $\le$ 2 intervalos & 503 logo após 2 intervalos & \Sim \\
    \addlinespace
    \cen{SE1} & chave inválida ou revogada: recusada, nada persistido & recusadas as duas; 0 linhas no central & \Sim \\
    \cen{SE2} & XSS e \emph{clickjacking} bloqueados & \emph{script} injetado e \texttt{<iframe>} bloqueados & \Sim \\
    \cen{SE3} & \emph{probe} para \emph{loopback} ou metadados recusada & 14 de 14 endereços recusados & \Sim \\
    \addlinespace
    \cen{IO1} & agente noutra linguagem só com o contrato & agente k6 em JavaScript; 0 perdas e 0 duplicados & \Sim \\
    \bottomrule
  \end{tabularx}
  \par\smallskip\footnotesize \Sim{} cumpre-se; \Parc{} cumpre-se em parte.
\end{table}

\section{Plano de evolução}
\label{sec:evolucao}

\subsection{Escalar para 10\,000+ hosts}
\porescrever{backlog E1--E8 com gatilhos medíveis (plano §8).}

\subsection{Decisões rejeitadas e adiadas}
\porescrever{plano §9.}

\section{Discussão}
\label{sec:discussao}

\subsection{Trade-offs observados}
\porescrever{o que as medições confirmaram ou contrariaram face à análise prévia.}

\subsection{Ameaças à validade}
\porescrever{uma só máquina; carga sintética; perfil de logs hipotético; métricas de código como proxies.}

\section{Conclusões}
\label{sec:conclusoes}

\porescrever{síntese, cumprimento dos objetivos e trabalho futuro.}


% ---------------------------------------------------- referências ---------
\clearpage
\phantomsection
\addcontentsline{toc}{section}{\refname}
\bibliographystyle{IEEEtran-pt}
\bibliography{referencias}

% -------------------------------------------------------- anexos ----------
\clearpage
\appendix
\renewcommand{\thesection}{\Alph{section}}
\section{Reprodução das medições}
\label{anx:reproducao}

\porescrever{comandos para reproduzir lint, testes, carga e E2E a partir do monitor-platform.}

\section{Registo de decisões arquiteturais}
\label{anx:adrs}

Cada decisão está registada num ADR em \texttt{docs/adr/} do repositório \texttt{monitor-platform}, com
o contexto e a evidência, a decisão, a teoria que a fundamenta, as alternativas rejeitadas, as
consequências e a forma de verificação (secção~\ref{sec:enquadramento}). A Tabela~\ref{tab:adrs} lista os
ADR; o texto completo de cada um está em \url{\Repositorio/tree/main/docs/adr}.

\begin{table}[htbp]
  \centering\small
  \caption{Decisões arquiteturais registadas}
  \label{tab:adrs}
  \begin{tabularx}{\linewidth}{lYlL{0.25\linewidth}}
    \toprule
    ADR & Decisão & Estado & QA \\
    \midrule
    0001 & Registar as decisões arquiteturais em ADRs & Aceite & Modifiability \\
    0002 & Repositório \texttt{monitor-platform} para CI partilhado e artefactos transversais & Aceite & Modifiability \\
    0003 & Contrato de ingestão em gRPC, num repositório \texttt{monitor-api} versionado & Aceite & Modifiability, Interoperability \\
    0004 & Medir primeiro: métricas Prometheus num listener interno e logs estruturados & Aceite & Monitorability \\
    0005 & Escritas em lote na ingestão e na resolução de alertas & Aceite & Scalability \\
    0006 & Controlo de admissão na ingestão, com ligações reservadas para o resto do central & Aceite & Scalability, Availability \\
    0007 & Central com várias instâncias: liderança por \emph{job}, limites partilhados e balanceador & Aceite & Scalability, Availability, Security \\
    0008 & Métricas de alerta como Strategy, com um registo como Simple Factory & Aceite & Modifiability \\
    0009 & \emph{Buffer} do remote com limite e descarte por severidade & Aceite & Availability, Performance \\
    0010 & Estado do \emph{buffer} reportado ao central; \texttt{/health/live} e \texttt{/health/ready} & Aceite & Monitorability, Availability \\
    0011 & O remote retoma a entrega sem reinício & Aceite & Availability \\
    0012 & Higiene dos repositórios e testes ponta a ponta em Compose & Aceite & Testability, Availability \\
    0013 & \emph{Probes} de saúde escaláveis: só os \emph{hosts} devidos, espalhados no tempo & Aceite & Availability, Performance \\
    0014 & \emph{Handlers} testáveis com interfaces mínimas; relógio e fuso horário injetados & Aceite & Testability \\
    0015 & Uma fonte única por repositório para os níveis de log e o endereço do central & Aceite & Modifiability \\
    0016 & Endurecimento por \emph{checklist} (ASVS 5.0 N1): cabeçalhos, CSP, Tailwind local e guarda SSRF & Aceite & Security \\
    0017 & Com 30 dias de dados, consultas do \emph{dashboard} com custo limitado e um \emph{bulkhead} no \emph{pool}; particionar só se a retenção o exigir & Aceite & Scalability, Performance \\
    0018 & PostgreSQL como base de dados única do central, para a telemetria e para o resto & Aceite (retroativo) & Scalability, Availability, Modifiability \\
    0019 & Estilo arquitetural: agentes \emph{store-and-forward} e um central monolítico em camadas, replicado sobre uma base de dados partilhada & Aceite (retroativo) & todos \\
    0020 & \emph{Store-and-forward} em SQLite no remote, com entrega \emph{at-least-once} e recibos de deduplicação no central & Aceite (retroativo) & Availability \\
    0021 & O SDK escreve diretamente no \emph{buffer} SQLite do remote, cujo esquema pertence ao remote & Aceite (retroativo) & Availability, Modifiability \\
    0022 & \emph{Outbox} transacional para os emails dos alertas & Aceite (retroativo) & Availability \\
    0023 & Chaves dos \emph{hosts} aleatórias de 256 bits, guardadas como \emph{digest} HMAC-SHA256 & Aceite (retroativo) & Security, Performance \\
    \bottomrule
  \end{tabularx}
\end{table}

Os ADR marcados como retroativos registam decisões que já estavam no código quando o trabalho começou, com a
justificação e a evidência que as medições entretanto deram (ADR-0001).


\end{document}
